% claim.tex -- AI4Math claim of record (AI-generated, 2026-09-30; instantiated
% from harness/templates/claim/claim.tex per SOP 08b).
% Machine-readable theorem anchors are the "% LEAN:" lines; scripts/paper-lint.sh
% and the claim audit check them. All Lean listings are verbatim, byte-identical
% contiguous excerpts of frozen artifacts of tasks/20260930-mossad42-quad:
%   statements  : formalized/03-half4-statements.lean (104 lines, sha256 b142670a...)
%   proof bodies: final/03-half4-proved.lean         (183 lines, sha256 548f3ef5...)
% Line ranges used: statement 30-103 (defs 34-50, 53-70, 77-81, 84 and signatures
% 87-88, 91-92, 95-96, 99-100, 103-104); the proof is a CONTIGUOUS EXCERPT.
% Verification: audit/listing-verify.txt (audit/inspect-tex.py --verify).
%
% SCOPE FENCE (must be preserved in any redistribution): the Lean layer below
% proves structural theorems about a pure finite-state integer recursion. The
% identification of that recursion with the constrained domino-tiling count is a
% combinatorial semantics bridge that is NOT proved here and is stated only as a
% conjecture. The definitions never mention tilings, dominoes or counting.
\documentclass[11pt]{article}

\usepackage[T1]{fontenc}
\usepackage[utf8]{inputenc}
\usepackage{amsmath,amssymb,amsthm}
\usepackage{graphicx}
\usepackage{hyperref}
\usepackage{xurl}
\setlength{\emergencystretch}{3em}
\input{lstlean}
\lstset{language=lean}

% --- local rendering patch (2026-09-30) ---
% tectonic/XeTeX does not apply the listings literate table to multi-byte
% characters, so the Unicode set used by the listings of THIS note is mapped to
% math glyphs as active characters via the standard \lccode/\lowercase idiom
% (ASCII-only source, no extra packages). Restricted to symbols outside xeCJK's
% own punctuation table: 00B7, 2013, 2014, 230A and 230B also occur in the
% frozen sources but are special punctuation that xeCJK sets up at load time --
% making them active first aborts the run with "Control sequence ... already
% defined" (observed 2026-09-27); those route through xeCJK + SimSun instead.
% 2026-09-30 修订（渲染面缺陷修复）：本清单此前只覆盖初稿用到的码点，后续
% 正文修订引入 ≠ ∏ ∈ ⇒ α ↦ ⟺ ①②③④ 而清单未同步 —— XeTeX 只发
% "Missing character" 警告并把字形**静默丢弃**（三道机械门都查不出）。现按
% claims/scan_listing_unicode.py 的 listing 内非 ASCII 码点盘点补齐；出包前须
% 以该盘点器无 MISSING 且编译日志 "Missing character" 计数为 0 复核。
% Rendering only: listing source bytes are untouched, so the byte-exactness
% checks are unaffected, and the \ifdefined guard leaves the pdflatex path
% identical to the template.
\ifdefined\XeTeXversion
  {\lccode`\~="2022 \lowercase{\global\catcode`~=\active \global\def~{\ensuremath{\bullet}}}}
  {\lccode`\~="2115 \lowercase{\global\catcode`~=\active \global\def~{\ensuremath{\mathbb{N}}}}}
  {\lccode`\~="2124 \lowercase{\global\catcode`~=\active \global\def~{\ensuremath{\mathbb{Z}}}}}
  {\lccode`\~="2190 \lowercase{\global\catcode`~=\active \global\def~{\ensuremath{\leftarrow}}}}
  {\lccode`\~="2192 \lowercase{\global\catcode`~=\active \global\def~{\ensuremath{\rightarrow}}}}
  {\lccode`\~="2194 \lowercase{\global\catcode`~=\active \global\def~{\ensuremath{\leftrightarrow}}}}
  {\lccode`\~="2200 \lowercase{\global\catcode`~=\active \global\def~{\ensuremath{\forall}}}}
  {\lccode`\~="2203 \lowercase{\global\catcode`~=\active \global\def~{\ensuremath{\exists}}}}
  {\lccode`\~="2227 \lowercase{\global\catcode`~=\active \global\def~{\ensuremath{\wedge}}}}
  {\lccode`\~="2228 \lowercase{\global\catcode`~=\active \global\def~{\ensuremath{\vee}}}}
  {\lccode`\~="2261 \lowercase{\global\catcode`~=\active \global\def~{\ensuremath{\equiv}}}}
  {\lccode`\~="2264 \lowercase{\global\catcode`~=\active \global\def~{\ensuremath{\leq}}}}
  {\lccode`\~="2265 \lowercase{\global\catcode`~=\active \global\def~{\ensuremath{\geq}}}}
  {\lccode`\~="22A2 \lowercase{\global\catcode`~=\active \global\def~{\ensuremath{\vdash}}}}
  {\lccode`\~="27E8 \lowercase{\global\catcode`~=\active \global\def~{\ensuremath{\langle}}}}
  {\lccode`\~="27E9 \lowercase{\global\catcode`~=\active \global\def~{\ensuremath{\rangle}}}}
  {\lccode`\~="00AC \lowercase{\global\catcode`~=\active \global\def~{\ensuremath{\neg}}}}
  {\lccode`\~="2212 \lowercase{\global\catcode`~=\active \global\def~{\ensuremath{-}}}}
  {\lccode`\~="2223 \lowercase{\global\catcode`~=\active \global\def~{\ensuremath{\mid}}}}
  {\lccode`\~="03B1 \lowercase{\global\catcode`~=\active \global\def~{\ensuremath{\alpha}}}}
  {\lccode`\~="21D2 \lowercase{\global\catcode`~=\active \global\def~{\ensuremath{\Rightarrow}}}}
  {\lccode`\~="2208 \lowercase{\global\catcode`~=\active \global\def~{\ensuremath{\in}}}}
  {\lccode`\~="2260 \lowercase{\global\catcode`~=\active \global\def~{\ensuremath{\neq}}}}
\fi
\ifdefined\XeTeXversion
  \usepackage{xeCJK}
  \setCJKmainfont{SimSun}
\fi

\newtheorem{theorem}{Theorem}

\title{Non-vanishing support for the half-horizontal-domino constrained tiling\\
count sequence on the $4\times n$ board\\
{\large (Lean 4 machine-checked claim of record)}}
\author{Aurora0134\thanks{Contact: via the GitHub account Aurora0134; ORCID
placeholder --- to be filled in by the author before any upload. No personal
information is carried in this note.}}
\date{\today}

\begin{document}
\maketitle

\begin{abstract}
We record a machine-checked result about a pure finite-state integer recursion: the
sequence $a_3(n) = D_4(n,0,n)$ built from a column-mask transfer table has non-vanishing
support exactly at the even indices, that is $a_3(n) \neq 0 \iff n \bmod 2 = 0$, together
with two state-level parity lemmas that every reachable state has an even mask popcount and
an even accumulated horizontal-tile count. All statements are verified by the Lean~4 kernel
under the pinned mathlib revision, with no \texttt{sorry} and axioms confined to
\{propext, \texttt{Quot.sound}, \texttt{Classical.choice}\}. Novelty: search found no
occupying record. This note is a priority claim of record, not a full paper.
\end{abstract}

\section{Introduction}
This note records a formalized result produced by an AI4Math automated pipeline, in which
propositions and proofs are generated by a language model and adjudicated solely by the
Lean~4 kernel. The result concerns a finite-state integer recursion: a transfer table on
column masks generates a two-parameter state count, from which the one-parameter sequence
$a_3$ is read off; the theorems are a parity invariant for reachable states and the exact
non-vanishing support of $a_3$. The note is a priority claim of record, deposited for
timestamping; the corresponding narrative paper has been deliberately deferred by the author
(see the card), so no full exposition is attempted here. Verification strength is stated in
Section~\ref{sec:verif}: strict kernel compilation, zero \texttt{sorry}, standard axiom
whitelist. The definitions never mention tilings, dominoes or counting, so the kernel layer
is a pure integer recursion. No literature survey is included---that expectation belongs to
arXiv moderation and does not apply to a Zenodo deposit.

\section{Statement}

We first fix notation. The recursion is built from a finite transfer table
\texttt{steps4} of 29 triples $(occ,next,j)$, where $occ$ is the incoming column mask,
$next$ the outgoing column mask and $j$ the number of horizontal tiles protruding from that
column; the table is generated by exhaustive enumeration of the four-row column-filling
recursion rather than transcribed by hand. Writing $\mathrm{pc}$ for the popcount of a mask
in the range $0..15$, and $\mathbf{0}$ for the empty mask, define the state count
\[
D_4(0,mask,h) = \begin{cases} 1, & mask = 0 \wedge h = 0,\\ 0, & \text{otherwise},\end{cases}
\qquad
D_4(n+1,mask,h) = \!\!\sum_{\substack{(occ,next,j) \in \texttt{steps4}\\ next = mask,\; j \le h}}\!\! D_4(n,occ,h-j),
\]
and the one-parameter sequence
\[
a_3(n) = D_4(n,\mathbf{0},n).
\]
The guard $j \le h$ is not removable: subtraction on $\mathbb{N}$ truncates, which would
pollute the count. The following five facts are formalized. The first two are state-level
parity lemmas, and the last is the main result.

% LEAN: tasks/20260930-mossad42-quad :: D4_pc_even  grade=完全证明
\begin{theorem}[reachable masks have even popcount; support lemma]
\label{thm:t3a}
For all $n, mask, h \in \mathbb{N}$, if $D_4(n,mask,h) \neq 0$ then $2 \mid \mathrm{pc}(mask)$.
\end{theorem}

% LEAN: tasks/20260930-mossad42-quad :: D4_h_even  grade=完全证明
\begin{theorem}[reachable states have even accumulated count; support lemma]
\label{thm:t3b}
For all $n, mask, h \in \mathbb{N}$, if $D_4(n,mask,h) \neq 0$ then $2 \mid h$.
\end{theorem}

% LEAN: tasks/20260930-mossad42-quad :: a3_odd_eq_zero  grade=完全证明
\begin{theorem}[vanishing at odd indices]
\label{thm:t3c}
For every $n \in \mathbb{N}$, if $n \bmod 2 = 1$ then $a_3(n) = 0$.
\end{theorem}

% LEAN: tasks/20260930-mossad42-quad :: a3_even_ne_zero  grade=完全证明
\begin{theorem}[non-vanishing at even indices]
\label{thm:t3d}
For every $k \in \mathbb{N}$, $\;a_3(2k) \neq 0$.
\end{theorem}

% LEAN: tasks/20260930-mossad42-quad :: a3_ne_zero_iff  grade=完全证明
\begin{theorem}[support characterization; main result]
\label{thm:t3e}
For every $n \in \mathbb{N}$, $\;a_3(n) \neq 0 \iff n \bmod 2 = 0$.
\end{theorem}

\begin{lstlisting}[caption={Frozen formal statement (Lean 4, verbatim from \texttt{formalized/03-half4-statements.lean}, lines 30--103), with the five proof-body placeholder lines omitted (see the note below the listing). sha256 \texttt{b142670a}...).},label={lst:stmt}]
import Mathlib

/-- **m=4 的列填充转移表**：`(occ, next, j)` 表示「入掩码 → 出掩码，该列伸出的水平砖数」。
穷举自 4 行的列填充递归（`phase0/audit-lean-steps.py` 生成，非手抄）。 -/
def steps4 : List (ℕ × ℕ × ℕ) :=
  [(0, 0, 0), (0, 12, 2), (0, 9, 2), (0, 3, 2), (0, 15, 4),
   (1, 8, 1), (1, 2, 1), (1, 14, 3),
   (2, 1, 1), (2, 13, 3),
   (3, 0, 0), (3, 12, 2),
   (4, 8, 1), (4, 11, 3),
   (5, 10, 2),
   (6, 9, 2),
   (7, 8, 1),
   (8, 4, 1), (8, 1, 1), (8, 7, 3),
   (9, 0, 0), (9, 6, 2),
   (10, 5, 2),
   (11, 4, 1),
   (12, 0, 0), (12, 3, 2),
   (13, 2, 1),
   (14, 1, 1),
   (15, 0, 0)]

/-- **掩码位数 `pc`（0..15 上的 popcount）**：显式列举（避免依赖额外位运算 API）。 -/
def pc : ℕ → ℕ
  | 0 => 0
  | 1 => 1
  | 2 => 1
  | 3 => 2
  | 4 => 1
  | 5 => 2
  | 6 => 2
  | 7 => 3
  | 8 => 1
  | 9 => 2
  | 10 => 2
  | 11 => 3
  | 12 => 2
  | 13 => 3
  | 14 => 3
  | 15 => 4
  | _ => 0

/-- **T3 的有限状态整数递推 `D4`**：
`D4 n mask h` = 填满前 n 列、出掩码 = mask、累计水平砖数 = h 的走法数。

注意 `j ≤ h` 保护不可省：ℕ 减法截断会污染计数（Phase 0 实测偏离，见
`phase0/audit-invariant-form.txt` 结论节）。 -/
def D4 : ℕ → ℕ → ℕ → ℕ
  | 0, mask, h => if mask = 0 ∧ h = 0 then 1 else 0
  | n + 1, mask, h =>
      (steps4.filter (fun t => t.2.1 = mask ∧ t.2.2 ≤ h)).map
        (fun t => D4 n t.1 (h - t.2.2)) |>.sum

/-- **T3 主序列 a3**：4×n 的「恰半」（累计水平砖 = n）走法数 = `D4 n 0 n`。 -/
def a3 (n : ℕ) : ℕ := D4 n 0 n

/-- **T3-支撑引理（I2）**：可达状态的掩码位数为偶（`pc mask ≡ 4·n ≡ 0 (mod 2)`）。 -/
theorem D4_pc_even (n mask h : ℕ) (hh : D4 n mask h ≠ 0) : 2 ∣ pc mask := by

/-- **T3-支撑引理（I1）**：可达状态的累计水平砖数恒为偶。 -/
theorem D4_h_even (n mask h : ℕ) (hh : D4 n mask h ≠ 0) : 2 ∣ h := by

/-- **T3-Z（零方向）**：`n` 为奇数时主序列取值为 0。 -/
theorem a3_odd_eq_zero (n : ℕ) (hn : n % 2 = 1) : a3 n = 0 := by

/-- **T3-N（非零方向）**：`n` 为偶数时主序列取值非零（4×2 块拼接构造）。 -/
theorem a3_even_ne_zero (k : ℕ) : a3 (2 * k) ≠ 0 := by

/-- **T3-I（支撑刻画，主定理）**：`a3 n` 非零当且仅当 `n` 为偶数。 -/
theorem a3_ne_zero_iff (n : ℕ) : a3 n ≠ 0 ↔ n % 2 = 0 := by
\end{lstlisting}

\noindent This is lines 30--103 of the frozen statement snapshot with exactly five lines
removed: the five proof-body placeholder lines (one per theorem). That single, stated elision
is the deposit convention (see the sibling claims): the frozen file's bodies are the
formalization department's placeholders, and reproducing an unfinished-proof marker in this
note would be rejected by the deposit's static lint. Everything else---doc-comments including
their internal notes, definitions, binders, signatures and the \texttt{:= by} keyword---is
verbatim and in order. \texttt{audit/inspect-tex.py} checks this mechanically: it asserts that
every remaining line occurs in the source in order, and that the only lines missing from the
block are those placeholder bodies. The real, sorry-free proofs follow in excerpt.

\section{Proof}
The complete machine proof is 183 lines long and therefore exceeds a comfortable page; the
listing below is a \emph{contiguous excerpt}, lines 86--183 of
\texttt{final/03-half4-proved.lean}, which is the entire proof development: it opens with the
proof-auxiliary section, contains the five private helper lemmas
(\texttt{map\_sum\_ne\_zero\_of\_mem}, \texttt{exists\_of\_sum\_ne\_zero},
\texttt{steps4\_guard}, \texttt{D4\_inv}, \texttt{D4\_step2}) that are additions on the proof
side and not part of the frozen statement, and then carries all five theorems with their
complete proof bodies. The material left out is lines 1--85 of the file: the header comment
(1--28), the \texttt{import} line, and the re-listing of the frozen definitions (30--84),
which are reproduced verbatim in the statement listing above. The full text is in the deposit
under \texttt{proofs/03-half4-proved.lean}. A prose reconstruction is not an obligation of
this note: narrative proofs belong to a full paper, which has been deferred.

\begin{lstlisting}[caption={Verbatim excerpt, lines 86--183 of \texttt{final/03-half4-proved.lean} (183 lines, sha256 \texttt{548f3ef5}...); the full text is in the deposit's \texttt{proofs/}.},label={lst:proof}]
/-! ## 证明辅助（军政部 T3 添加，不改动上方冻结 statement） -/

/-- 列表 `map f` 之和非零 ⇒ 某个被映射元素的值非零。 -/
private lemma map_sum_ne_zero_of_mem {α : Type*} (f : α → ℕ) (l : List α) {a : α}
    (ha : a ∈ l) (hfa : f a ≠ 0) : (l.map f).sum ≠ 0 := by
  intro h0
  exact hfa (List.sum_eq_zero_iff.mp h0 (f a) (List.mem_map.mpr ⟨a, ha, rfl⟩))

/-- `map f` 之和非零 ⇒ 存在一项其值非零。 -/
private lemma exists_of_sum_ne_zero {α : Type*} (f : α → ℕ) (l : List α)
    (h : (l.map f).sum ≠ 0) : ∃ a ∈ l, f a ≠ 0 := by
  by_contra hc
  simp only [not_exists, not_and, not_not] at hc
  exact h (List.sum_eq_zero_iff.mpr fun x hx => by
    obtain ⟨a, ha, rfl⟩ := List.mem_map.mp hx
    exact hc a ha)

/-- `steps4` 的转移恒等式（`decide` 穷举 29 条全验）：列格数偶 + 出掩码位数 = 伸出砖数。 -/
private lemma steps4_guard {occ mask j : ℕ} (ht : (occ, mask, j) ∈ steps4) :
    2 ∣ (pc occ + j) ∧ pc mask = j := by
  have hall : steps4.all
      (fun t => decide (2 ∣ (pc t.1 + t.2.2) ∧ pc t.2.1 = t.2.2)) = true := by decide
  exact of_decide_eq_true (List.all_eq_true.mp hall (occ, mask, j) ht)

/-- **I1 ∧ I2（互归纳的合并不变量）**：可达状态的累计水平砖数与掩码位数均为偶。 -/
private lemma D4_inv (n : ℕ) :
    ∀ mask h : ℕ, D4 n mask h ≠ 0 → 2 ∣ h ∧ 2 ∣ pc mask := by
  induction n with
  | zero =>
    intro mask h hh
    by_cases hc : mask = 0 ∧ h = 0
    · obtain ⟨rfl, rfl⟩ := hc; decide
    · exact absurd (by simp [D4, hc]) hh
  | succ n ih =>
    intro mask h hh
    have hh' : ((steps4.filter (fun t => t.2.1 = mask ∧ t.2.2 ≤ h)).map
        (fun t => D4 n t.1 (h - t.2.2))).sum ≠ 0 := hh
    obtain ⟨t, htmem, htne⟩ := exists_of_sum_ne_zero (fun t => D4 n t.1 (h - t.2.2))
      (steps4.filter (fun t => t.2.1 = mask ∧ t.2.2 ≤ h)) hh'
    obtain ⟨htsteps, htp⟩ := List.mem_filter.mp htmem
    obtain ⟨hocc, hjle⟩ := of_decide_eq_true htp
    obtain ⟨ihh, ihpc⟩ := ih t.1 (h - t.2.2) htne
    obtain ⟨hg1, hg2⟩ := steps4_guard htsteps
    have hdj : 2 ∣ t.2.2 := by
      have h := Nat.dvd_sub hg1 ihpc
      rwa [Nat.add_sub_cancel_left] at h
    have hdh : 2 ∣ h := by
      have h := Nat.dvd_add ihh hdj
      rwa [Nat.sub_add_cancel hjle] at h
    exact ⟨hdh, by rw [← hocc, hg2]; exact hdj⟩

/-- **T3-N 的构造步**：两步环 `(0,12,2)` → `(12,0,0)` 把非零走法升 2 列、加 2 砖。 -/
private lemma D4_step2 (n h : ℕ) (hh : D4 n 0 h ≠ 0) : D4 (n + 2) 0 (h + 2) ≠ 0 := by
  have h1 : D4 (n + 1) 12 (h + 2) ≠ 0 := by
    refine map_sum_ne_zero_of_mem (f := fun t => D4 n t.1 (h + 2 - t.2.2))
      (l := steps4.filter (fun t => t.2.1 = 12 ∧ t.2.2 ≤ h + 2)) (a := (0, 12, 2)) ?_ ?_
    · apply List.mem_filter.mpr
      exact ⟨by decide, by simp⟩
    · simpa using hh
  refine map_sum_ne_zero_of_mem (f := fun t => D4 (n + 1) t.1 (h + 2 - t.2.2))
    (l := steps4.filter (fun t => t.2.1 = 0 ∧ t.2.2 ≤ h + 2)) (a := (12, 0, 0)) ?_ ?_
  · apply List.mem_filter.mpr
    exact ⟨by decide, by simp⟩
  · simpa using h1

/-- **T3-支撑引理（I2）**：可达状态的掩码位数为偶（`pc mask ≡ 4·n ≡ 0 (mod 2)`）。 -/
theorem D4_pc_even (n mask h : ℕ) (hh : D4 n mask h ≠ 0) : 2 ∣ pc mask :=
  (D4_inv n mask h hh).2

/-- **T3-支撑引理（I1）**：可达状态的累计水平砖数恒为偶。 -/
theorem D4_h_even (n mask h : ℕ) (hh : D4 n mask h ≠ 0) : 2 ∣ h :=
  (D4_inv n mask h hh).1

/-- **T3-Z（零方向）**：`n` 为奇数时主序列取值为 0。 -/
theorem a3_odd_eq_zero (n : ℕ) (hn : n % 2 = 1) : a3 n = 0 := by
  by_contra hc
  have h2 : 2 ∣ n := D4_h_even n 0 n (by simpa [a3] using hc)
  rw [Nat.dvd_iff_mod_eq_zero] at h2
  omega

/-- **T3-N（非零方向）**：`n` 为偶数时主序列取值非零（4×2 块拼接构造）。 -/
theorem a3_even_ne_zero (k : ℕ) : a3 (2 * k) ≠ 0 := by
  induction k with
  | zero => decide
  | succ k ih =>
    exact D4_step2 (2 * k) (2 * k) (by simpa [a3] using ih)

/-- **T3-I（支撑刻画，主定理）**：`a3 n` 非零当且仅当 `n` 为偶数。 -/
theorem a3_ne_zero_iff (n : ℕ) : a3 n ≠ 0 ↔ n % 2 = 0 := by
  constructor
  · intro h
    by_contra hodd
    have : n % 2 = 1 := by omega
    exact h (a3_odd_eq_zero n this)
  · intro h
    have hn : n = 2 * (n / 2) := by omega
    rw [hn]
    exact a3_even_ne_zero (n / 2)
\end{lstlisting}

\section{Verification evidence}
\label{sec:verif}
\begin{itemize}
\item Toolchain: Lean 4 \texttt{v4.34.0}; mathlib4 \texttt{v4.34.0}, revision \texttt{5ed29652}.
\item Statement integrity: \texttt{formalized/03-half4-statements.lean}, 104 lines,
sha256 \texttt{b142670a02c923ff2d90e416957938e222762f139a67d5f64b593ce9064ed74b}.
Final artifact \texttt{final/03-half4-proved.lean}, 183 lines,
sha256 \texttt{548f3ef528b432c4c0cbfb95de7993b246d1cc782f9ec2042952af68db30036a}.
Every listing above is byte-identical to its source (checked by \texttt{audit/inspect-tex.py},
which asserts that each listed line occurs in its source in order and that the only omitted
lines are the statement file's placeholder proof bodies).
\item Kernel check: strict compilation exits 0 with no \texttt{sorry} and no \texttt{admit}.
\texttt{\#print axioms} output, verbatim:
\begin{verbatim}
'D4_pc_even' depends on axioms: [propext, Classical.choice, Quot.sound]
'D4_h_even' depends on axioms: [propext, Classical.choice, Quot.sound]
'a3_odd_eq_zero' depends on axioms: [propext, Classical.choice, Quot.sound]
'a3_even_ne_zero' depends on axioms: [propext, Classical.choice, Quot.sound]
'a3_ne_zero_iff' depends on axioms: [propext, Classical.choice, Quot.sound]
\end{verbatim}
Every line is a subset of \{propext, \texttt{Quot.sound}, \texttt{Classical.choice}\};
\texttt{sorryAx} does not occur.
\item Audit chain: statement freeze (sha256), byte-wise symmetric-difference comparison against
the frozen snapshot, strict-mode compilation, and the axiom whitelist check. The comparison
records 9 of 9 frozen declarations byte-identical and notes that the proof artifact adds five
auxiliary lemmas, which is allowed. Originals are in the deposit under \texttt{audit/}.
\end{itemize}

\section{Novelty evidence}
Under the pipeline's C9 novelty discipline the verdict for this object is \emph{no occupying
record found}. The full query log, including every zero-hit query, is in the deposit at
\texttt{audit/c9-record.md}; it is reproduced verbatim from the screening batch and the
present batch's supplementary search.

Summary of channels. OEIS: the sequence, its compressed subcolumn, its even-index variant and
its first difference all returned zero hits. Short prefixes of the compressed subcolumn
matched two entries (A220379 and A292752), but the match breaks at the fifth term; these are
short-window coincidences, do not constitute an occupying record, and are deliberately not
cited here as related literature. Channel availability was established by positive controls
before any zero-hit was read: A001835 and A005178 (the unconstrained $3\times n$ and
$4\times n$ domino-tiling counts) matched. Literature: four channels were queried (arXiv API,
Crossref, Mathematics Stack Exchange, zbMATH Open); there were no direct hits on this
constrained-selection object. Library layer: no hits in mathlib, compfiles, or the scanned
external repositories.

Two honest limitations are recorded and must not be dropped in any restatement.
First, the OpenAlex channel was \emph{unreachable} from the machine used for this search
(Cloudflare-fronted hosts timed out; DNS resolution was verified correct against a public
resolver, so this is not a DNS-poisoning artifact). Crossref was used as a substitute
DOI-level index. Coverage is therefore \emph{lower} than in the screening batch, and this is
a channel gap, not strengthened evidence.
Second, a zero hit is \emph{not} novelty: the strongest claim licensed here is that no record
of this constrained-selection object was found in the channels queried.
No moment or expectation novelty is claimed: mixed moments $E[V^a H^b]$ for general $m\times n$
boards are already occupied by \emph{Statistics of Domino Tilings on a Rectangular Board},
Fibonacci Quarterly (2019), doi:10.1080/00150517.2019.12427625~\cite{fq2019}.
The $3\times n$ and $4\times n$ objects of this family are $m$-dimensional variants of one
another and are not used as independent novelty evidence for each other; the variant partner
of the present $4\times n$ track is the $3\times n$ half-horizontal track, which is not cited
here as support.

\section{Scope and limitations}
The kernel-level results concern the pure finite-state integer recursion defined above. They
do \emph{not} establish that $a_3(n)$ equals the number of domino tilings of the $4\times n$
board in which horizontal dominoes make up exactly half of all tiles. That identification is
a combinatorial semantics bridge; it is stated here only as a conjecture, is not part of the
theorem layer, and is supported only by numerical probes and external sequence-database
evidence. The definitions themselves never mention tilings, dominoes or counting---the state
is a column mask and the recursion involves only masks and integer counts---so what is
formalized is a pure integer recursion, and the combinatorial bridge is still \emph{not} in
the theorem layer. Accordingly, the results are described throughout as structural theorems
about a recursion, never as a proved statement about tiling counts.

The definition layer of this track was changed, under a user-approved rule-override event
recorded as O-2, from a pure sequence recurrence to a finite-state integer recursion (a
dynamic-programming layer) with $D_4\colon \mathbb{N}\to\mathbb{N}\to\mathbb{N}\to\mathbb{N}$
and $a_3(n) := D_4(n,0,n)$. This is a change of the object being formalized, not a weakening:
the grading of all five theorems is \emph{complete proof} (0 \texttt{sorry}, axiom whitelist),
and no extra hypothesis was added. The absence of a constant-coefficient (C-finite) recurrence
for $a_3$ must be read as \emph{none found at order $\le 60$}, not as \emph{does not exist};
the order-4 P-recursive relation reported during exploration was not brought into the frozen
statement and is not claimed here, so this note claims neither a recurrence nor a closed form.

This note contains no literature survey; it is a deposit record, not a survey paper.

\section*{Data availability}
Lean sources, verification and novelty-search evidence are deposited on Zenodo
(DOI: \texttt{RESERVED-DOI}) and mirrored as a public snapshot repository. The note is
released under CC BY 4.0 and the code under MIT.

\section*{Use of generative AI}
\begin{itemize}
\item \emph{AI-performed stages}: proposition generation, natural-language-to-Lean
formalization, proof search, and drafting of this text.
\item \emph{Model, node, budget}: pipeline node \texttt{anthropic/a6api-main/deepseek-v4.1-flash}
(wire-id; resolved at level L2 of the pipeline's model-detection ladder, snapshot taken
2026-09-30 09:35). Source-task budget: sampling 1 of 8 allowed runs on this track (10 of 32
across four parallel tracks); judgement-class LLM calls 6 of 16 across the batch; Lean
formal compilations 6 of 12 on this track (19 of 48 across the batch); LaTeX 0. The numbers are
copied from the source task's \texttt{budget.log}.
\item \emph{Final adjudication}: all statements and proofs reported here were checked by the
Lean~4 kernel. Every proof compiles with no \texttt{sorry}, and \texttt{\#print axioms}
reports only \{propext, \texttt{Quot.sound}, \texttt{Classical.choice}\}. Statement text was
frozen before proving and compared byte-wise afterwards.
\item The author reviewed the material and is responsible for all content. AI systems are
not listed as authors.
\end{itemize}

\section*{Author contributions}
The author determined the research question and the scope fences, verified the semantics of
the statements, and is responsible for signing off and for any deposit or release. The
automated pipeline (LLM-driven) generated the proposition, the formalization, the proof
search and the draft text. No external release is performed by the pipeline.

\begin{thebibliography}{9}
% Only works actually cited; each was checked to exist before entry. Verification
% criterion is the returned title (not merely the HTTP status), following the
% 2026-09-28 erratum lesson where a DOI tail was mis-resolved to a different article.
\bibitem{lean4}
L.~de Moura and S.~Ullrich.
\newblock The Lean 4 theorem prover and programming language.
\newblock \emph{CADE-28}, 2021. \url{https://doi.org/10.1007/978-3-030-79876-5_37}

\bibitem{mathlib}
The mathlib Community.
\newblock The Lean mathematical library.
\newblock \emph{CPP 2020}, 2020. \url{https://doi.org/10.1145/3372885.3373824}

\bibitem{fq2019}
\emph{Statistics of Domino Tilings on a Rectangular Board}.
\newblock Fibonacci Quarterly, 2019.
\newblock \url{https://doi.org/10.1080/00150517.2019.12427625}
\end{thebibliography}

\end{document}
